\documentclass[11pt]{article}
\usepackage{amsmath,amssymb,amsthm}
\usepackage[margin=1in]{geometry}
\newtheorem{theorem}{Theorem}
\newtheorem{proposition}[theorem]{Proposition}
\newtheorem{lemma}[theorem]{Lemma}
\newtheorem{corollary}[theorem]{Corollary}
\theoremstyle{definition}\newtheorem{definition}[theorem]{Definition}
\newcommand{\gen}[1]{\langle #1\rangle}
\newcommand{\Z}{\mathrm Z}
\title{Normal subgroups forced by an end intersection}
\author{Working mathematical note}
\date{}
\begin{document}
\maketitle

This note concerns abstract finite groups only.  It extends the internal
note \emph{The intersection property in rank three for simple groups} by
retaining an additional generator that centralizes the first two generators.
The dihedral intersection argument is borrowed from that note and proved
again here, including its exceptional case of order four.  The extension
identifies a normal solvable obstruction in the ambient group; it does not
require either end subgroup to be simple.  No historical novelty claim is
made.  There is no numerical calculation, classification claim, formalization,
or credit for an Original case in this document.

\begin{definition}
For a labelled generating family $(r_i)_{i\in E}$ of a group, write
$G_A=\gen{r_i:i\in A}$ for $A\subseteq E$, with $G_\varnothing=1$.
The family has the \emph{intersection property} if
\[
                 G_A\cap G_B=G_{A\cap B}
                 \qquad(A,B\subseteq E).
\]
A family of nonidentity involutions in a specified order is a string family
if nonadjacent generators commute.  A string family with the intersection
property is a string C-group representation.
\end{definition}

We use the following condition on a finite group:
\begin{equation}\label{eq:radical}
 \text{$G$ has no nonidentity solvable normal subgroup.}
\end{equation}
This is often expressed by saying that its solvable radical is trivial.
Only the displayed condition is used; no structure theorem about finite
groups is required.

\begin{theorem}\label{thm:ends}
Let $G$ be a finite group satisfying \eqref{eq:radical}, and suppose that
\[
 G=\gen{a,b,c,d},\qquad a^2=b^2=c^2=d^2=1,
 \qquad a,b,c,d\ne1,
\]
\[
                         [a,c]=[a,d]=[b,d]=1.
\]
Then $(\gen{a,b,c};a,b,c)$ and $(\gen{b,c,d};b,c,d)$ both have the
intersection property.  In particular, the assertion holds when $G$ is a
finite nonabelian simple group.
\end{theorem}

We give the elementary ingredients before proving the theorem.  The
subgroups in its conclusion can be proper and need not be simple.

\begin{lemma}\label{lem:dihedral}
Let $x,y$ be distinct nonidentity involutions in a finite group, put
$D=\gen{x,y}$ and $R=\gen{xy}$, and let $B\le D$ contain $y$.
Then $R$ is cyclic of index two in $D$, every element of $D\setminus R$
is an involution, and
\[
                 B=(B\cap R)\rtimes\gen y.
\]
Every subgroup of $R$ is normal in $D$.  If $|B|>4$, then $B\cap R$ is
the unique cyclic subgroup of index two in $B$.  If $|B|=4$, then $B$ is
a Klein four group and $B\cap R=\gen z$, where $z$ is the central
involution of $D$ lying in $R$.
\end{lemma}
\begin{proof}
Set $u=xy$ and $m=|u|$.  The identity $yuy=u^{-1}$ puts every element of
$D$ in the form $u^i y^\epsilon$, where $0\le i<m$ and
$\epsilon\in\{0,1\}$.  The involution $y$ cannot lie in $R$: otherwise
$D=R$ would be cyclic, whereas a finite cyclic group has at most one
nonidentity involution and $x,y$ are distinct.  The two normal-form cosets
are therefore disjoint, and $|D|=2m$.  Also
$(u^iy)^2=1$.  The quotient map $D\to D/R$ restricts surjectively to
$B$, since $y\in B\setminus R$, proving the asserted decomposition.

Any subgroup of the cyclic group $R$ is preserved by conjugation by $u$
and by inversion, hence by $D=\gen{u,y}$.  In particular, it is normal
in $D$; no normality of $B$ in $D$ is being assumed.

Write $k=|B\cap R|$, so that $|B|=2k$.  If $k>2$, a generator of any
cyclic index-two subgroup of $B$ has order $k>2$.  All elements outside
$B\cap R$ are involutions.  This generator must therefore lie in
$B\cap R$, and the equality of orders proves uniqueness.

If $k=2$, let $z$ generate $B\cap R$.  The involutions $z,y$ commute
because $y$ inverts $z$, so $B=\gen{z,y}$ is a Klein four group.
As $R$ is cyclic, $m$ is even and $z=u^{m/2}$.  This element is fixed
by inversion and commutes with $u$, hence is central in $D$.
\end{proof}

The uniqueness assertion deliberately excludes a Klein four group.  Its
three cyclic subgroups of index two are different.  Generalized quaternion
groups cannot occur in the lemma: the displayed splitting has an
involution outside the cyclic subgroup, and all elements of that outside
coset have order two.  The degenerate dihedral group of order four is
included explicitly; a subgroup $B$ of order two is simply $\gen y$.

\begin{proposition}\label{prop:obstruction}
Let $G=\gen{a,b,c,d}$ be finite.  Suppose that $a,b,c$ are distinct
nonidentity involutions, $[a,c]=1$, and $d$ centralizes $\gen{a,b}$.
The element $d$ need not be an involution in this proposition.  Put
\[
 P=\gen{a,b},\qquad Q=\gen{b,c},\qquad I=P\cap Q.
\]
If $I\ne\gen b$, then $G$ has either a nonidentity cyclic normal
subgroup or a nonidentity normal subgroup generated by two involutions.
In either event, $G$ has a nonidentity solvable normal subgroup.
\end{proposition}
\begin{proof}
Let $R_1=\gen{ab}$ and $R_2=\gen{bc}$, and set $J_i=I\cap R_i$.
The common generator $b$ is outside both rotation subgroups.  By
Lemma~\ref{lem:dihedral}, each $J_i$ is cyclic of index two in $I$ and
$I=J_i\rtimes\gen b$.

Suppose first that $|I|>4$.  Uniqueness of the cyclic index-two subgroup
gives $J_1=J_2=:J$.  It is nontrivial and normal in both $P$ and $Q$:
each $J_i$ is a subgroup of its corresponding cyclic rotation group.
The element $d$ centralizes $P$ and therefore centralizes $J$.
Consequently $J$ is normal in $\gen{P,Q,d}=G$.  This produces the first
asserted obstruction.

Suppose next that $|I|=4$.  The lemma gives central involutions
\[
                z_1\in\Z(P),\qquad z_2\in\Z(Q),
                \qquad I=\gen{b,z_1}=\gen{b,z_2}.
\]
If $z_1=z_2$, this nonidentity element centralizes $P,Q$ and $d$.
Its cyclic subgroup is central, hence normal in $G$.

It remains to consider $z_1\ne z_2$.  Neither equals $b$, so the three
nonidentity elements of the Klein four group $I$ are $b,z_1,z_2$, and
\begin{equation}\label{eq:bproduct}
                              b=z_1z_2.
\end{equation}
Define
\[
                   A=\gen{a,z_2},\qquad C=\gen{c,z_1}.
\]
Every generator of $A$ commutes with every generator of $C$: the four
relations are $[a,c]=1$, $[a,z_1]=1$, $[z_2,c]=1$ and
$[z_2,z_1]=1$.  The middle two use centrality in $P$ and $Q$, respectively;
the last uses the abelian group $I$.  Thus $[A,C]=1$.
Both factors lie in $H=\gen{a,b,c}$, and their product contains $a,c$
and $b=z_1z_2$.  Hence $H=AC$, and $A$ is normal in $H$.

Moreover, $d$ commutes with $a,b$ and $z_1\in P$.  By
\eqref{eq:bproduct}, it also commutes with $z_2=bz_1$.  It therefore
centralizes $A$.  This proves $A\lhd\gen{H,d}=G$.  The subgroup $A$
is nontrivial, since it contains $a\ne1$, and is generated by two
involutions.  No distinctness or noncommutation of these two auxiliary
involutions is needed.

For completeness, every group generated by two involutions is solvable:
their product generates a cyclic normal subgroup with quotient of order
at most two.  Its derived subgroup lies in that cyclic subgroup, and its
second derived subgroup is trivial.  This also covers a cyclic subgroup
of order two and a Klein four group.  Thus the subgroup $A$ just obtained
is a solvable normal subgroup.

Finally $|I|$ is even because $b\in I$ has order two.  If neither of
the preceding cases occurs, $|I|=2$ and $I=\gen b$, contrary to the
assumption of the proposition.  The cases are exhaustive.
\end{proof}

\begin{lemma}\label{lem:degeneracies}
Under the hypotheses of Theorem~\ref{thm:ends}, the four generators are
pairwise distinct and no adjacent pair commutes.  Consequently each
adjacent product has order at least three.
\end{lemma}
\begin{proof}
The centre of $G$ is trivial, since any nonidentity central element
generates a nonidentity cyclic normal subgroup, forbidden by
\eqref{eq:radical}.  Each possible generator equality would supply a
nonidentity central generator, as follows:
\[
\begin{array}{c|cccccc}
\text{equality}&a=b&a=c&a=d&b=c&b=d&c=d\\\hline
\text{central generator}&a&d&a&b&a&c.
\end{array}
\]
Each entry follows by substituting the indicated equality into
$[a,c]=[a,d]=[b,d]=1$.  This excludes all six equalities.

If $[a,b]=1$, the generator $a$ commutes with all four generators and
is central.  If $[c,d]=1$, the same applies to $d$.  If $[b,c]=1$,
both $c$ and $d$ centralize $P=\gen{a,b}$, so $P\lhd G$.
It is a nonidentity group generated by two involutions and hence solvable,
again contradicting \eqref{eq:radical}.  No adjacent pair commutes.
For involutions $x,y$, the equality $(xy)^2=1$ is equivalent to
$xy=yx$; together with pairwise distinctness this proves the last claim.
\end{proof}

\begin{lemma}\label{lem:rankthree}
Let $a,b,c$ be distinct nonidentity involutions with $ac=ca$.  If
\begin{equation}\label{eq:smallintersection}
                 \gen{a,b}\cap\gen{b,c}=\gen b,
\end{equation}
then $(\gen{a,b,c};a,b,c)$ has the full intersection property.
\end{lemma}
\begin{proof}
Write $P=\gen{a,b}$, $Q=\gen{b,c}$ and $V=\gen{a,c}$.
The group $V$ is a Klein four group.  The equality
\eqref{eq:smallintersection} implies $c\notin P$ and $a\notin Q$:
otherwise the nonidentity involution $c$ or $a$, distinct from $b$,
would belong to $P\cap Q=\gen b$.  A subgroup of $V$ containing $a$
is either $\gen a$ or $V$.  Since $c\notin P$, it follows that
\[
               P\cap V=\gen a,\qquad Q\cap V=\gen c.
\]
Together with \eqref{eq:smallintersection}, these are all intersections
of distinct two-generator subgroups.  In particular $b\notin V$, since
$b\in P$ and $P\cap V=\gen a$ with $b\ne a$.

A singleton subgroup therefore meets the subgroup on the other two
generators trivially.  Distinct singleton subgroups meet trivially as
their nonidentity involutions are different.  All other subset pairs are
nested, involve the empty or full subset, or are the pairs of
two-element subsets already treated.  These cases exhaust the subsets
of $\{a,b,c\}$, giving the stated full property.
\end{proof}

\begin{proof}[Proof of Theorem~\ref{thm:ends}]
The hypotheses of Proposition~\ref{prop:obstruction} hold for the ordered
quadruple $(a,b,c,d)$, by Lemma~\ref{lem:degeneracies} and the string
relations.  Condition~\eqref{eq:radical} rules out both conclusions of
that proposition.  Hence
\[
                         \gen{a,b}\cap\gen{b,c}=\gen b.
\]
Lemma~\ref{lem:rankthree} proves the intersection property in the first
end subgroup.

Apply exactly the same argument to the reversed quadruple $(d,c,b,a)$.
Its first and third entries commute because $[d,b]=1$, and its fourth
entry centralizes its first two because $[a,d]=[a,c]=1$.  We obtain
\[
                         \gen{d,c}\cap\gen{c,b}=\gen c,
\]
and the intersection property in the second end subgroup follows.

Finally, a nonabelian simple group is not solvable.  Indeed, its derived
subgroup is nontrivial, since it is nonabelian, and is normal, hence is
the whole group.  Its derived series cannot reach the identity.  A
nonidentity solvable normal subgroup of a simple group would be the
whole group, an impossibility.  Thus finite nonabelian simple groups
satisfy \eqref{eq:radical}.
\end{proof}

\section*{A conditional consequence in rank four}

The following elementary criterion records exactly what remains once the
two end subgroups are known to have the intersection property.  Its proof
is included to distinguish a statement about the ends from the full
rank-four assertion.

\begin{lemma}\label{lem:maximal}
For a generating family indexed by a finite set $E$, suppose that every
subfamily indexed by $E\setminus\{i\}$ has the intersection property,
and that the subgroups indexed by maximal proper subsets satisfy
\[
 G_{E\setminus\{i\}}\cap G_{E\setminus\{j\}}
                    =G_{E\setminus\{i,j\}}\qquad(i\ne j).
\]
Then the full family has the intersection property.
\end{lemma}
\begin{proof}
Only proper subsets $A,B$ need consideration.  Choose maximal proper
index sets $X\supseteq A$ and $Y\supseteq B$.  If $X=Y$, use the
intersection property in $G_X$.  Otherwise the assumed maximal
intersection gives $G_A\cap G_B\le G_{X\cap Y}$.
The intersection property within $G_X$ and $G_Y$ gives
\[
 G_A\cap G_{X\cap Y}=G_{A\cap Y},\qquad
 G_B\cap G_{X\cap Y}=G_{B\cap X}.
\]
Both groups on the right have index sets contained in $X$.
Intersecting them within $G_X$ therefore yields $G_{A\cap B}$,
as required.
\end{proof}

\begin{proposition}\label{prop:gluing}
Suppose that $a,b,c,d$ are distinct nonidentity involutions generating
$G$, with $[a,c]=[a,d]=[b,d]=1$.  If both end subgroups have the
intersection property and
\begin{equation}\label{eq:endintersection}
                \gen{a,b,c}\cap\gen{b,c,d}=\gen{b,c},
\end{equation}
then $(G;a,b,c,d)$ has the full intersection property.
\end{proposition}
\begin{proof}
Write $P=\gen{a,b}$, $M=\gen{b,c}$, $T=\gen{c,d}$ and
\[
 H_0=\gen{b,c,d},\quad H_1=\gen{a,c,d},\quad
 H_2=\gen{a,b,d},\quad H_3=\gen{a,b,c}.
\]
The end properties and \eqref{eq:endintersection} imply $a\notin H_0$
and $d\notin H_3$.  For example, $a\in H_0$ would put $a$ in
$H_0\cap H_3=M$, contrary to the intersection property in $H_3$.
Since $a$ centralizes $T$ and $d$ centralizes $P$, it follows that
\[
                  H_1=\gen a\times T,\qquad
                  H_2=P\times\gen d.
\]
Both displayed families have the intersection property: each factor
does, and intersections of subgroups generated by subsets of the given
generators are computed coordinatewise in a direct product.
Thus all four maximal subfamilies have that property.

For clarity we check all six maximal intersections.  The first is the
hypothesis $H_0\cap H_3=M$.  The direct product decompositions and
$a\notin H_0$, $d\notin H_3$ give
\[
                  H_0\cap H_1=T,\qquad H_2\cap H_3=P.
\]
For $x=a^\epsilon t\in H_1\cap H_3$, one has
$t\in T\cap H_3=T\cap M=\gen c$, where the last equality is in
the end subgroup $H_0$.  Hence $H_1\cap H_3=\gen{a,c}$.
The symmetric calculation gives $H_0\cap H_2=\gen{b,d}$.

Finally $P\cap T=1$.  Indeed, this intersection lies in $H_0\cap H_3=M$;
the end properties give $P\cap M=\gen b$ and $T\cap M=\gen c$,
whose intersection is trivial.  If
\[
                 x=a^\epsilon t=p d^\eta\in H_1\cap H_2,
                 \quad p\in P,\ t\in T,
\]
then $a^{-\epsilon}p=t d^{-\eta}\in P\cap T=1$.
Therefore $p=a^\epsilon$, $t=d^\eta$, and
$H_1\cap H_2=\gen{a,d}$.  All six intersections have their required
forms.  Lemma~\ref{lem:maximal} completes the proof.
\end{proof}

\begin{corollary}\label{cor:onecondition}
Under the hypotheses of Theorem~\ref{thm:ends}, the rank-four family
$(G;a,b,c,d)$ has the intersection property if and only if
\eqref{eq:endintersection} holds.  In particular, for a finite nonabelian
simple group the full rank-four question reduces to this one end
intersection, once the stated membership, generation, involution and
string hypotheses are established.
\end{corollary}
\begin{proof}
Necessity is part of the definition.  For sufficiency use
Theorem~\ref{thm:ends} and Proposition~\ref{prop:gluing}.
\end{proof}

\section*{Why the ambient hypothesis matters}

It would be incorrect to discard \eqref{eq:radical} or replace it by
an assumption merely about the number of generators.  One can see the
order-four obstruction even with distinct generators and all adjacent
products of order greater than two.
Let $S_3=\gen{s,t}$ with $s,t$ distinct transpositions, and put
\[
 G=S_3\times S_3,\quad
 a=(s,1),\ b=(t,t),\ c=(1,s),\ d=(1,t).
\]
These are distinct nonidentity involutions satisfying the string
relations.  They generate $G$, because $bd=(t,1)$ and the two factors
are generated by their transpositions.  The adjacent product orders
are six, six and three.  Nevertheless
\[
 \gen{a,b}=S_3\times\gen t,
 \qquad \gen{b,c}=\gen t\times S_3,
\]
so their intersection is $\gen t\times\gen t$, of order four, rather
than $\gen{(t,t)}$, of order two.  For the first equality, $(ab)^3=(1,t)$;
the second follows similarly from $(bc)^3=(t,1)$.  The normal dihedral
factor supplied by Proposition~\ref{prop:obstruction} is the first copy
of $S_3$.  This example also shows why the end subgroup must not be
silently treated as simple.

\section*{Conditions for a later application}

An application to a specific marking must separately identify its finite
ambient group satisfying \eqref{eq:radical}, establish that all four
marks belong to that group and generate it, and establish the nonidentity,
involution and nonadjacent commutation assertions.  The resulting end
intersection must be established in the same group and faithful frame.
Only then does Corollary~\ref{cor:onecondition} yield the full rank-four
intersection property.  This note supplies none of those numerical or
case-specific premises.

In particular, the abstract theorem can replace a separate assertion of
the end rank-three intersection properties.  It cannot supply a complete
middle-dihedral journal, a claimed dihedral order, a faithful frame,
membership or generation, tensor congruences, or the measured end
intersection.  Those remain independent premises of any certificate
that uses them.  No missing historical arithmetic body is recreated here.
\end{document}
